\section*{Resumen}

El presente proyecto tiene como objetivo desarrollar una herramienta tecnológica
que permita a los productores de kéfir artesanal optimizar el tiempo de recolección
del kéfir y determinar si su fermentación se llevó a cabo de manera óptima, con el
fin de garantizar la calidad del producto. Para ello, se desarrollarán cuatro módulos.
El primer módulo se encargará de obtener los datos del proceso de fermentación
mediante sensores. Estos datos serán procesados en el segundo módulo, donde se
realizará un análisis utilizando un modelo basado en reglas de lógica difusa para
evaluar el estado de la fermentación. Luego, el tercer módulo almacenará tanto los
datos obtenidos de los sensores como los resultados del análisis en una base de
datos. Finalmente, el cuarto módulo permitirá visualizar toda esta información a
través de una página web. Este sistema proporcionará a los productores de kéfir
artesanal la posibilidad de analizar sus fermentaciones basándose en datos
estructurados, en lugar de depender exclusivamente de conocimientos empíricos.
\newline
\newline
% Palabras clave
\noindent \textbf{Palabras clave:} kéfir, fermentación, sensores, análisis de datos, lógica difusa, base de datos, página web.
